\section{Three fixed simples and an anti-invariant module}
\label{sec:three-fixed}
The following argument uses only semisimple module categories and the
ordinary fusion rules. In particular, its conclusions are general.

\begin{proposition}[Anti-invariant module representation]
\label{prop:anti-module}
Suppose a commutative fusion ring is categorified, an invertible simple
$g$ has order two, and exactly $f$ simples $F_1,\ldots,F_f$ are fixed
by $g$. There are integral $f\times f$ matrices $T_i$ satisfying
\[
 T_0=I,\qquad T_{i^*}=T_i^{\mathsf t},\qquad
 T_iT_j=\sum_kN_{ij}^kT_k,
\]
\[
 |(T_i)_{ab}|\leq N_{iF_a}^{F_b},\qquad
 (T_i)_{ab}\equiv N_{iF_a}^{F_b}\pmod2.
\]
Each $T_i$ is normal, and its spectrum is contained in that of $N_i$.
\end{proposition}
\begin{proof}
A fixed simple gives a rank-one module category over the pointed
subcategory generated by $g$, trivializing its associator. Thus
$A=1\oplus g$ is a connected separable algebra. Each fixed simple has
two right $A$-module structures, since $H^2(C_2,\C^\times)=0$; a free
orbit gives one simple module. Twisting by the algebra automorphism
that is $-1$ on $g$ exchanges the two structures over each $F_a$ and
fixes the free-orbit modules. This involution commutes with the left
category action. On its minus eigenspace, use the basis
$F_a^+-F_a^-$. The two nonnegative action coefficients have sum
$N_{iF_a}^{F_b}$, so their difference has the asserted bounds and
parity. Module reciprocity gives transpose duality. Commutativity then
gives normality. The squarefree characteristic polynomial of $N_i$
annihilates $b_i$ in the semisimple complexified fusion ring, and hence
annihilates $T_i$.
\end{proof}

\begin{theorem}
\label{thm:three-fixed}
None of
\[
 \begin{gathered}
 R7\_2\_116,\ R7\_2\_261,\ R7\_3\_364,\ R7\_3\_938,\\
 R7\_3\_970,\ R7\_4\_2644,\ R7\_4\_2683
 \end{gathered}
\]
is categorifiable by a complex fusion category.
\end{theorem}
\begin{proof}
All indices below are source indices, starting with $X_0=1$.

For $R7\_2\_116$ and $R7\_2\_261$, take respectively
$(g,X)=(X_1,X_2)$ and $(X_3,X_1)$, with fixed ordering
$(X_4,X_5,X_6)$. The bound matrix for $T_X$ is
\[
 B=\begin{pmatrix}0&1&1\\1&0&1\\1&1&2\end{pmatrix}.
\]
The source spectra lie respectively in $\{0,\pm1,5\}$ and
$\{0,\pm1,\pm i,5\}$. The maximum absolute row sum is at most four,
excluding eigenvalue five. Normality implies operator norm at most
one. Its first row nevertheless contains two entries of absolute
value one, a contradiction.

For $R7\_3\_364$ and $R7\_3\_938$, take respectively $g=X_1,X_3$,
again with fixed ordering $(X_4,X_5,X_6)$. The self-dual simple $X_4$
has bound matrix and squarefree source polynomial
\[
 B=\begin{pmatrix}3&1&2\\1&1&2\\2&2&2\end{pmatrix},\qquad
 p(x)=x(x^2-8x-8)(x^2-2x-2).
\]
Write its symmetric anti-invariant matrix as
\[
 T=\begin{pmatrix}a&b&c\\b&d&e\\c&e&f\end{pmatrix},
 \qquad a,b,d\text{ odd},\quad c,e,f\text{ even}.
\]
Its integral characteristic polynomial has degree three. Thus its
spectrum is three zeros, $\{0,4\pm2\sqrt6\}$, or
$\{0,1\pm\sqrt3\}$. The first option contradicts symmetry and odd
entries. The second has $\operatorname{tr}(T^2)=80$, whereas the
bounds give $\operatorname{tr}(T^2)\leq32$. In the last option the
trace, trace-square, and determinant are $2,8,0$. The trace-square
and parity force $c=e=0$, $|a|=|b|=|d|=1$, and $|f|=2$.
The determinant forces $ad=b^2=1$, so $a=d$. Consequently
$a+d+f\in\{-4,0,4\}$, contradicting trace two.

For $R7\_3\_970$, take $g=X_3$ and fixed ordering
$(X_4,X_5,X_6)$. The bound matrix for $T_1$ is
\[
 \begin{pmatrix}0&0&0\\0&1&2\\0&2&1\end{pmatrix},
\]
and its source spectrum lies in $\{0,\pm1,\pm i,9\}$. The row-sum
bound three excludes nine. Normality and parity therefore give
$T_1=T_2=D=\operatorname{diag}(0,\varepsilon,\eta)$, where
$\varepsilon,\eta=\pm1$. The respective bounds on $T_5,T_6$ are
\[
 \begin{pmatrix}0&2&0\\2&2&0\\0&0&0\end{pmatrix},\qquad
 \begin{pmatrix}0&0&2\\0&0&0\\2&0&2\end{pmatrix}.
\]
Commutation with $D$ kills both off-diagonal pairs. The remaining
diagonal supports are disjoint, so $T_5T_6=0$. The source product
$X_5X_6=2X_1+2X_2$ instead gives $T_5T_6=4D\ne0$.

For $R7\_4\_2644$, take $g=X_3$, fixed ordering
$(X_4,X_5,X_6)$, and test simple $X_1$. Its bounds and polynomial are
\[
 \begin{pmatrix}1&2&2\\2&1&2\\2&2&3\end{pmatrix},\qquad
 (x-1)(x+1)(x^2+1)(x^2-14x+1).
\]
The row-sum bound seven excludes $7+4\sqrt3$, and rationality of the
characteristic polynomial excludes its conjugate. Normality and
parity imply $T_1=T_2=D$, diagonal with entries $\pm1$. The first
column of $T_5$ is zero by the source bounds. Thus $(T_4T_5)_{00}=0$,
whereas $X_4X_5=2X_1+2X_2$ forces this entry to be $\pm4$.

Finally, for $R7\_4\_2683$, take $g=X_5$ and fixed ordering
$(X_1,X_2,X_6)$. The same bounds and polynomial, now for $X_3$, give
$T_3=T_4=D$, diagonal with entries $\pm1$. The first row of $T_1$ is
zero, while the source product $X_1^2=2X_3+2X_4$ forces
$(T_1^2)_{00}=4D_{00}=\pm4$. This completes all seven contradictions.
\end{proof}

The accompanying source verifier reconstructs the involutions, fixed
simples, duality, bound matrices, spectral polynomials, and displayed
products directly from the pinned census. The proof does not require
a search through associators or sphericalization signs.
